\documentclass[10pt,a4paper]{amsart}
\usepackage[margin=19mm]{geometry}
\usepackage{amsmath,amssymb,mathtools}
\usepackage[scr=rsfs,cal=euler]{mathalfa}
\usepackage{xfrac}
\usepackage[unicode,hidelinks,bookmarks=false]{hyperref}
\newcommand{\ie}{i.e.\ }
\newcommand{\cf}{cf.\ }
\newcommand{\resp}{respectively\ }
\newcommand{\Subsection}{\subsection}
\providecommand{\nobreakdash}{\nobreak\hbox{-}}
\setlength{\emergencystretch}{2em}
\multlinegap0pt
\usepackage{bm}
\usepackage{bbm}
\usepackage{dsfont}
\usepackage{xspace}
\usepackage{cancel}
\usepackage{mathtools}
\usepackage{amsthm}
\makeatletter
\@ifundefined{thm}{\newtheorem{thm}{Thm}}{}
\@ifundefined{lemma}{\newtheorem{lemma}[thm]{Lemma}}{}
\makeatother
\makeatletter
\newcommand{\Tt}{\mathcal{T}}
\newcommand{\acc}{\mathrm{acc}}
\newcommand{\al}{{\alpha}}
\expandafter\ifx\csname example\endcsname\relax
\newenvironment{example}{\begin{example_no_diamond} } {\hfill$\er$ \end{example_no_diamond}}
\fi
\expandafter\ifx\csname itemlist\endcsname\relax
\newenvironment{itemlist}
   { \begin{list} {$\bullet$}
         { \setlength{\topsep}{.5ex}  \setlength{\itemsep}{.5ex} \setlength{\leftmargin}{2.5ex} } }
\fi
\newcommand{\la}{{\lambda}}
\newcommand{\ov}{\overline}
\expandafter\ifx\csname rmk\endcsname\relax
\newenvironment{rmk}{\begin{rmk_no_diamond} } {\hfill$\er$ \end{rmk_no_diamond}}
\fi
\makeatother
\title[Quadrilateral mutations and symplectic embeddings]{Quadrilateral mutations and symplectic embeddings}
\author{Nicki Magill}
\date{Source: arXiv:2606.12729v1 (CC BY 4.0)}
\begin{document}
\maketitle

\noindent\emph{Source and licence.} Quadrilateral mutations and symplectic embeddings. Version arXiv:2606.12729v1 (CC BY 4.0), distributed under the Creative Commons Attribution 4.0 International licence, as recorded in the arXiv licence metadata for this exact version and shown on the abstract page https://arxiv.org/abs/2606.12729v1. Full rights notice: licenses/arxiv-2606.12729-CC-BY-4.0.txt.\par\smallskip

\noindent\textit{Editorial scope.} Counted unit: the Lemma whose source label is \texttt{lem:yTailXVisible} in the article (source0.tex lines 3194--3202, source label \texttt{lem:yTailXVisible}), delivered together with the article's own complete original proof of it (source0.tex lines 3204--3281, one \texttt{proof} environment of 78 lines). No mathematics is written, repaired or re-derived: statement and proof are the article's own text line for line. Required context retained uncounted: lem:fullFillingAccPt (lines 3001--3004); lem:pqAlOrder2 (lines 3124--3127); lem:affineDegenerationLimitH (lines 3166--3173). Every definition, display and label of those lines is retained unchanged. Omitted from this package and not redistributed: 1--3000, 3005--3123, 3128--3165, 3174--3193, 3203--3203, 3282--3537. Nothing inside a retained excerpt is omitted or altered: every retained body is delivered exactly as the article's lines are written, so no omission receipt of any kind accompanies this package. Citations: the retained excerpts carry no citation command at all, so no bibliography entry of the article is transcribed and no reference is redistributed. Reference adaptation: none was needed and none was made; every reference inside the delivered unit resolves inside it, so no unresolved reference or citation is produced. Printed slips kept literal: no printed word, symbol, hypothesis or number of the counted statement or of its proof was altered. Layout and class: the article's own class is \texttt{amsart} with packages including etex, textcmds, amsmath, amssymb, amsfonts, amstext, verbatim, amsthm, mathrsfs, microtype, xy, lineno, xcolor, aliascnt; the wrapper is the lane's pinned \texttt{amsart} editorial wrapper at 10pt on A4, which restates the theorem-like environments the retained text needs and adds the article's own macro definitions that the retained text needs (\texttt{\textbackslash Tt}, \texttt{\textbackslash acc}, \texttt{\textbackslash al}, \texttt{\textbackslash la}, \texttt{\textbackslash ov}), with extra wrapper packages (bm, bbm, dsfont, xspace, cancel, mathtools). The delivered numbering is the unit's own. Assets: no retained span contains a figure environment, a graphics-inclusion command, a PDF-page inclusion, a table or a TikZ picture; no image file is copied into this package, the package declares no asset dependency and no image file is redistributed with it. Licence: version arXiv:2606.12729v1 of this article is distributed under the Creative Commons Attribution 4.0 International licence, as recorded in the arXiv licence metadata for this exact version and shown on the abstract page https://arxiv.org/abs/2606.12729v1. A full rights notice is delivered at licenses/arxiv-2606.12729-CC-BY-4.0.txt. Characters: every retained excerpt is pure ASCII, so the delivered engine, XeLaTeX with its default Latin Modern text font, reports no missing character.
\begin{lemma} \label{lem:fullFillingAccPt}
   Let $w_k=x^kw'$ (resp. $y^kw'$) where $w'$ is some word in $x,y,s,\ov{x},\ov{y},\ov{s}.$  Assume that $w_kQ_{\bullet,b}^{0}:=OX_kV_kY_k$ is well-defined for each $k$ for either $\bullet=P,H$. If the $\lim_{k \to \infty} w_kQ_{\bullet,b}^{0}$ is such that the limit of the affine length of $V_kY_k$ (resp. $X_kV_k$) goes to zero and 
   the limit of the slope of $X_kV_k$ (resp. $V_kY_k$) is irrational, then $c_{\bullet,b}(z_\infty)$ is unobstructed. 
\end{lemma}

\begin{lemma} \label{lem:pqAlOrder2}
    Let $\Tt=(\bm{x}_\la,\bm{x}_\mu,\bm{x}_\rho)$ be a correctly ordered triple for $\bullet=H,P$. Then, 
    \[ p_\la/q_\la<p_\mu/q_\mu<\acc_\bullet(\al_\mu)<p_\rho/q_\rho<\acc_\bullet(\al_\rho)\]
\end{lemma}

\begin{lemma}\label{lem:affineDegenerationLimitH}
 For $\bullet=H,P$, assume that
\[
        Q_k:=Q_{\bullet,b}(y^k\Tt)=OXV_kY_k
\]
is well-defined for all $k$, and that the affine length of the edge
$XV_k$ tends to zero. Then, in the case of $H_b$, we must have $b=\lim_{k \to \infty} \frac{m_k}{d_k},$ and in the case of $P_b,$ we must have $b=\lim_{k \to \infty} \frac{e_k}{f_k}$. 
\end{lemma}

\begin{lemma}\label{lem:yTailXVisible}
Let $\Tt$ be a correctly ordered recursive triple. 
Assume that
\[
        Q_k:=Q_{\bullet,b}(y^k\Tt)=OXV_kY_k
\]
is well-defined for all $k$ and $Q_k$ affinely degenerates to a triangle for $\bullet=H,P$. In the case of $H_b$, we assume $b \neq 1/3,$ and in the case of $P_b$ we assume $b \neq 1.$ Then $xQ_k$ is well-defined for all sufficiently large $k$.

\end{lemma}

\begin{proof}
As $Q_k$ degenerates to a triangle and $X$ is constant, we must have $|XV_k| \to 0.$
Define
\[
        y^k\Tt=(\bm{x}_{k-1},\bm{x}_k,\bm{x}_\rho).
\]
Let
\[
        A:=|OX|,
        \qquad
        B:=\lim_{k\to\infty}|OY_k|.
\]
By the affine triangular degeneration, the limiting triangle has vertices
\[
        O=(0,0),\qquad X=(A,0),\qquad Y_\infty=(0,B).
\]
Let
\[
        z_\infty:=\acc_\bullet(b).
\]
Then, by the proof of Lemma~\ref{lem:fullFillingAccPt}, for the limiting triangle,
\[
        \left\{\frac{A}{B},\frac{B}{A}\right\}
        =
        \left\{z_\infty,\frac{1}{z_\infty}\right\}.
\]
For each $k$, the nodal ray at $X$ of $Q_k$ is constant and $\vec n_X=(p_\rho-6q_\rho,\ q_\rho).$ Note, if $\ov{s}Q_k$ is defined for just some $k,$ this immediately implies that $p_\rho-6q_\rho<0$. 

Suppose, for contradiction, that $\ov{s}Q_k$ is well-defined for infinitely
many $k$. Passing to an infinite subsequence if necessary, the nodal ray at
$X$ intersects the side $OY_k$ for all $k$ in this subsequence. Taking the
limit, the same ray must intersect the segment $OY_\infty$. Hence, we get
\begin{equation} \label{eq:bound6prho}
        A \frac{q_\rho}{6q_\rho-p_\rho}\leq B \iff  \frac{p_\rho}{q_\rho}
        \leq
        6-\frac{A}{B}.
\end{equation}


On the other hand, since the triples $y^k\Tt$ are correctly ordered,
Lemma~\ref{lem:pqAlOrder2} gives
\[
        \acc_\bullet(\alpha_k)<\frac{p_\rho}{q_\rho}
\]
for every $k$. Taking $k\to\infty$ and using that $\alpha_k\to b$ by Lemma~\ref{lem:affineDegenerationLimitH}, we obtain
\[
        z_\infty\leq \frac{p_\rho}{q_\rho}.
\]

We now compare the two bounds. In the case of $H_b$ we have $b \neq 1/3$ and in the case of $P_b$ we have $b \neq 1,$ hence in either case, the accumulation point
satisfies
\[
        z_\infty>3+2\sqrt{2} \iff  z_\infty+\frac{1}{z_\infty}>6
\]

Using this and the fact that $A/B$ is either $z_\infty$ or $1/z_\infty,$ we conclude that 

\[
        6-\frac{A}{B}<z_\infty \leq p_\rho/q_\rho,
\]
but this contradicts \eqref{eq:bound6prho}.

Therefore $\ov{s}Q_k$ can be well-defined for only
finitely many $k$. 

If the nodal ray at \(X\) hit the vertex \(Y_k\) for infinitely many \(k\),
then passing to a subsequence and taking the limit would give
\[
        \frac{p_\rho}{q_\rho}=6-\frac{A}{B}.
\]
But we have shown
\[
        6-\frac{A}{B}<z_\infty\leq \frac{p_\rho}{q_\rho},
\]
so this equality is impossible. Hence the vertex-hit case occurs only
finitely many times. Therefore $xQ_k$ is
well-defined for all sufficiently large $k$.
\end{proof}

\end{document}
